\section{Step 25 Inspect the provided vector quantizer before switching}\label{fir-s25}

The vector path already contains the rounding, bounds, endpoint replacement, counting and byte stores just developed in the scalar path. Compare its fir labels with their scalar counterparts. The final vector sum must arrive in \texttt{t0}; \texttt{v0} retains the original count and \texttt{v1} accumulates clipped outputs. Reusing the established rule avoids introducing a different quantizer while changing the dot product.

\subsection{Expected observation}
The vector postprocessing advances \texttt{a0} by two, keeps the coefficient base, and repeats \texttt{fir\_window}. It applies one bias and arithmetic shift after each complete sum.

\Needspace{11\baselineskip}
\subsection{Check your understanding}
What must the vector reduction supply to the provided postprocessing?

\begin{itemize}
\item \textbf{A} One exact complete accumulator in \texttt{t0}.
\item \textbf{B} Only the low word of a pair sum.
\item \textbf{C} Eight separately rounded halfwords.
\end{itemize}
